immini[thm]{\subsection*{I. Áp suất khí theo mô hình động học phân tử}
\subsubsection{Tác dụng của một phân tử khí lên thành bình}
\begin{itemize}
	\item Xét một lượng khí gồm N phân tử chứa trong một bình lập phương có cạnh $\ell$, trong hệ tọa độ vuông góc Oxyz.
	\item Một phân tử có khối lượng m chuyển động thẳng đều song song với trục Ox với tốc độ v từ thành bình EFOH tới va chạm đàn hồi và trực diện với thành bình ABCD.
	\item Sau khi va chạm, phân tử chuyển động theo chiều ngược lại với tốc độ có cùng độ lớn v tới thành bình đối diện.
\end{itemize}}{\begin{tikzpicture}[scale=1,thick,>=stealth']
\tkzDefPoints{0/0/e, 0.75/0.5/f, 2.25/0.5/b, 1.5/0/a, 0/-1.5/h, 1.5/-1.5/d, 2.25/-1/c, 0.75/-1/o,-0.75/-2/z, 3/-1/x, 0.75/1.25/y, 0.6/-0.8/m, 1.2/-0.8/n}
\draw[fill=green] (0,0) -- (0.75,0.5) -- (2.25,0.5) -- (1.5,0) -- (0,0);
\draw[fill=green!30] (2.25,0.5) -- (2.25,-1) -- (1.5,-1.5) -- (1.5,0);
\draw[fill=green!5] (0,0) -- (0,-1.5) -- (1.5,-1.5) -- (1.5,0) -- (0,0);
\tkzDrawSegments[thick,->](o,z o,x o,y)
\tkzDrawSegments[](a,b e,f f,b e,a e,h h,d d,c c,b d,a)
\draw[->,red] (1.2,-0.8) -- (1.7,-0.8);
\draw[->,red] (0.6,-0.8) -- (0.1,-0.8);
\draw[fill=red,red] (m) circle(.8pt);
\draw[fill=red,red] (n) circle(.8pt);
\tkzLabelPoint[above left](a){$A$}
\tkzLabelPoint[above](b){$B$}
\tkzLabelPoint[left](e){$E$}
\tkzLabelPoint[left](f){$F$}
\tkzLabelPoint[above right](o){$O$}
\tkzLabelPoint[below](c){$C$}
\tkzLabelPoint[below](d){$D$}
\tkzLabelPoint[below](h){$H$}
\tkzLabelPoint[below](z){$z$}
\tkzLabelPoint[below](x){$x$}
\tkzLabelPoint[left](y){$y$}
\end{tikzpicture}
}
\Binhluan{\begin{itemize}
	\item Do độ lớn của vận tốc tức thời và tốc độ tức thời của một chuyển động có độ lớn bằng nhau nên động lượng của phân tử trước khi va chạm với thành bình có giá trị là $+mv$, sau khi va chạm với thành bình có giá trị là $-mv$. Độ biến thiên động lượng của phân tử do va chạm với thành bình ABCD có độ lớn là:
	\begin{align*}
		\left|\Delta \vec p \right|=\left|-mv-(+mv) \right|=|-2mv|=2mv
	\end{align*} 
	\item Coi chuyển động trước và sau khi va chạm với thành bình là chuyển động thẳng đều.
\end{itemize}
}{Chú ý về các kí hiệu: 
\begin{itemize}
	\item $\vec p$: động lượng
	\item p: áp suất
	\item $p_m$: áp suất của một phân tử khí.
	\end{itemize}}
\begin{itemize}
	\item Thời gian giữa hai va chạm liên tiếp của một phân tử lên thành bình ABCD: $\Delta t=\dfrac{2\ell}{v}$

	\item Lực do phân tử khí tác dụng lên thành bình có độ lớn: $\vec F.\Delta t=\Delta \vec p \Rightarrow F=\dfrac{\Delta p}{\Delta t}=\dfrac{mv^2}{\ell}$
	\item Áp suất do một phân tử khí tác dụng lên thành bình ABCD có giá trị là:
	\begin{align*}
	 p_m=\dfrac{F}{S}=\dfrac{F}{\ell^2}=\dfrac{mv^2}{\ell^3}=\dfrac{m}{V}.v^2
	 \end{align*}
	 với thể tích lượng khí $V=\ell^3$.
\end{itemize}
\subsubsection{Tác dụng của N phân tử khí lên thành bình}
\begin{itemize}
	\item Vì số phân tử N vô cùng lớn và các phân tử chuyển động hỗn loạn trong bình nên các hướng Ox, Oy và Oz là bình đẳng.
	\item Do đó, số phân tử chuyển động theo hướng Ox, từ mặt EFOH sang mặt ABCD và ngược lại để gật áp suất lên hai mặt này chỉ bằng $\dfrac{1}{3}$ số phân tử trong bình $\left(\dfrac{N}{3}\right)$.
	\item Vì các phân tử chuyển động hỗn loạn nên tốc độ của các phân tử không bằng nhau, do đó áp suất mỗi phân tử tác dụng lên thành bình cũng không bằng nhau.
	\item Trung bình mỗi phân tử tác dụng lên thành bình một áp suất $\overline{p}_m=\dfrac{m}{V}.\overline{v^2}$, trong đó: 
	\begin{align*}
		\overline{v^2}=\dfrac{v_1^2+v_2^2+...+v_N^2}{N}
	\end{align*}
	\item Đại lượng này được gọi là "trung bình của các bình phương tốc độ".
	\item Từ những lập luận trên ta có thể viết được công thức tính áp suất của khí trong bình tác dụng lên thành bình ABCD: 
	\begin{align*}
		p=\dfrac{N}{3}.\overline{p}_m=\dfrac{1}{3}.\dfrac{Nm}{V}.\overline{v^2}\quad (*)
	\end{align*}
	\item Trong đó $\mu= \dfrac{N}{V}$ là số phân tử có trong một đơn vị thể tích, gọi là \textit{mật độ phân tử}.
	\item Động năng trung bình của phân tử là $\overline{E}_{\text{đ}}=\dfrac{m\overline{v^2}}{2}$ và thay $\overline{E}_{\text{đ}}$ vào phương trình $(*)$ ta được: 
	\begin{align*}
		p=\dfrac{2}{3}\mu \overline{E}_{\text{đ}}
	\end{align*}
	\item Vì các mặt của hình lập phương là bình đẳng nên áp suất của chất khí tác dụng lên các mặt của bình đều như nhau.
\end{itemize}
\Binhluan{\subsection*{II. Mối quan hệ giữa động năng phân tử và nhiệt độ}
\begin{itemize}
	\item Dựa vào các công thức $pV=nRT$, $\mu=\dfrac{N}{V}$, $N=n.N_A$ và $p=\dfrac{2}{3}\mu \overline{E}_{\text{đ}}$
\begin{align*}
&\Rightarrow \overline{E}_{\text{đ}}=\dfrac{3p}{2\mu}=\dfrac{3}{2}.\dfrac{nRT}{\mu V}=\dfrac{3}{2}.RT.\dfrac{n}{\dfrac{N}{V}.V}= \dfrac{3}{2}.RT.\dfrac{n}{N}\\ 
&\Rightarrow \boder{\overline{E}_{\text{đ}}=\dfrac{3}{2}.\dfrac{R}{N_A}T}
\end{align*}
\item Vì R và $N_A$ đều là các hằng số xác định, nên ta có thể tính được hằng số k:
\begin{align*}
	k=\dfrac{R}{N_A}=\dfrac{8,31\ J/molK}{6,02.10^{23}\ \text{hạt/mol}}=1,38.10^{-23}\ J/K
\end{align*}
\item Hằng số k được gọi là hằng số Boltzmann.
\item Do đó động năng trung bình của phân tử được xác định bằng hệ thức: 
\begin{align*}
	\boder{\overline{E}_{\text{đ}}=\dfrac{3}{2}kT}
\end{align*}
\item Từ hệ thức trên ta thấy động năng trung bình của các phân tử tỉ lệ thuận với nhiệt độ tuyệt đối.
\end{itemize}
\newghichu{
$\bullet$ Các khí có bản chất khác nhau, khối lượng khác nhau nhưng nhiệt độ như nhau thì động năng trung bình của các phân tử bằng nhau.\\
$\bullet$ Động năng trung bình của phân tử khí càng lớn thì nhiệt độ của khí càng cao.\\
$\bullet$ Vì $\overline{E}_{\text{đ}}$ tỉ lệ thuận với T nên người ta có thể coi nhiệt độ tuyệt đối là số đo động năng trung bình của phân tử theo một đơn vị khác.
}
\begin{note}
	Căn bậc hai của $\overline{v^2}$ là $\sqrt{\overline{v^2}}$, độ lớn của đại lượng này không phải là trung bình cộng của các tốc độ phân tử, nghĩa là không phải là tốc độ trung bình của các phân tử. Nó được gọi là \textit{tốc độ căn quân phương} của phân tử.
\end{note}
}
{Hệ thức $\overline{E}_{\text{đ}}=\dfrac{3}{2}kT$ chỉ đúng cho khí lí tưởng, trong đó phân tử được coi là chất điểm, chuyển động của phân tử là chuyển động tịnh tiến.\\
Do đó, động năng trung bình trong hệ thức trên được gọi một cách đẩy đủ là động năng tịnh tiến trung bình.\\
Khi phân tử không được coi là chất điểm thì ngoài chuyển động tịnh tiến, phân tử còn có thể có chuyển động quay, chuyển động riêng của các nguyên tử trong phân tử, chuyển động riêng của các electron trong nguyên tử,... Khi đó cần thiết phải nêu rõ động năng trung bình trong biểu thức trên chỉ là động năng của chuyển động tịnh tiến của phân tử, gọi tắt là động năng tịnh tiến trung bình.
\begin{align*}
	\overline{E}_{\text{đ}}=\dfrac{m}{2}.\overline{v^2}
\end{align*}
} 
